
When a language model performs well on MMLU and poorly on HellaSwag, practitioners often attribute this asymmetry to the structure of the benchmarks themselves. A complementary question, less often posed, is what differs in the \textit{pre-training data} that produced this asymmetry. Data mixing decisions --- which domains to include, in what proportions, across how many tokens --- are among the most consequential and least understood choices in language model development. This paper attempts to measure which domains of The Pile \citep{Gao2020} drive which NLP benchmarks, and characterizes both what can be confirmed and what remains open.

The domain-benchmark specificity question has a structural answer in the literature: data composition matters in aggregate. DCLM \citep{Li2024} demonstrates that model-based filtering outperforms heuristic filtering (64\% vs. approximately 57\% MMLU at 7B scale). Domain mixing frameworks such as RegMix \citep{Liu2024} and DoReMi \citep{Xie2023} show that optimal domain weights improve aggregate validation loss. Topic Over Source \citep{Peng2025} establishes that topic-based mixing consistently outperforms source-based mixing. But these analyses report \textit{aggregate} effects across tasks. No study has produced a per-domain $\times$ per-benchmark coefficient matrix from a controlled within-family analysis. It is not known empirically whether Wikipedia exposure specifically improves MMLU more than HellaSwag, or whether Books exposure specifically improves commonsense reasoning more than factual recall.

The absence of this result is not due to lack of interest. Comparing model families introduces architecture and training confounds. Re-training with varied domain proportions is computationally prohibitive at scale. The essential requirement is within-family domain exposure variation that can serve as a covariate in a regression mapping domain composition to benchmark performance --- and this variation has not been characterized at the per-checkpoint level.

Our key observation is that this variation already exists in the Pythia model family \citep{Biderman2023}. Because The Pile stores documents in domain-clustered blocks rather than uniformly shuffling them, different training checkpoints of the same Pythia model have seen different cumulative proportions of Wikipedia, StackExchange, PubMed Abstracts, and other domains. Pythia's identity document index mapping (\texttt{doc\textbackslash \_idx.npy}, confirmed for 134M entries) allows exact reconstruction of which documents each checkpoint consumed, and therefore how much of each domain it had been exposed to at each point in training. This creates a natural panel dataset: 154 checkpoints across multiple model sizes, each with a distinct domain exposure profile, enabling regression analysis of domain-benchmark specificity without new training runs.

We operationalize this design through four sub-hypothesis experiments spanning domain content analysis (h-m1), exposure trajectory extraction (h-e1), Spearman correlation analysis (h-m2), and panel OLS regression (h-m3). The results are empirically coherent but incomplete. The two foundational experiments pass their gates convincingly. The two specificity experiments fail their gates due to a shared structural data limitation: Books3 domain exposure is zero across all 154 checkpoints in the 600,000-document first-shard sample. This is not a pipeline error but a consequence of The Pile's block structure, which concentrates Books3 documents in later shards not covered by the proof-of-concept lookup.

This paper makes the following contributions:

\begin{enumerate}
\item \textbf{Empirical grounding for domain-benchmark specificity analysis.} To our knowledge, this is the first direct quantitative demonstration that The Pile's domains are separable by automated cognitive task pattern proxies with effect sizes exceeding $\eta^2$ = 0.91. This confirms the empirical basis for the assumption underlying domain-differentiated data mixing theories. These measurements were conducted on domain-representative generated texts; $\eta^2$ values are upper bounds relative to real Pile documents.
\end{enumerate}

\begin{enumerate}
\item \textbf{Reusable infrastructure for within-family domain exposure analysis.} We extract and validate domain exposure trajectories from Pythia's exact dataloaders using the identity \texttt{doc\textbackslash \_idx} mapping, providing a proven pipeline (\texttt{compute\textbackslash \_domain\textbackslash \_exposure\textbackslash \_trajectories()}, \texttt{build\textbackslash \_domain\textbackslash \_lookup()}, \texttt{step\textbackslash \_to\textbackslash \_sample()}) for future domain attribution studies.
\end{enumerate}

\begin{enumerate}
\item \textbf{Identification of a structural data scope requirement.} Single-shard Pile samples systematically miss 6 of 22 domains (including Books3 and GitHub). The full 134M-document lookup across all 30 Pile shards is the minimum required dataset for domain-benchmark specificity analysis.
\end{enumerate}

\begin{enumerate}
\item \textbf{Preliminary directional evidence at 70M scale.} The observed reversal ($\rho$(Wikipedia, HellaSwag) > $\rho$(Wikipedia, MMLU) at 70M) is most consistently explained by floor effects at 70M scale and N = 10 checkpoint sampling artifact, but raises the open question of whether domain-benchmark specificity is a scale-emergent phenomenon.
\end{enumerate}

